
\documentclass{article}
\usepackage{colortbl}
\usepackage{makecell}
\usepackage{multirow}
\usepackage{supertabular}

\begin{document}

\newcounter{utterance}

\centering \large Interaction Transcript for game `cladder', experiment `full\_v1.5\_default', episode 4119 with llama{-}3.1{-}8b{-}instruct.
\vspace{24pt}

{ \footnotesize  \setcounter{utterance}{1}
\setlength{\tabcolsep}{0pt}
\begin{supertabular}{c@{$\;$}|p{.15\linewidth}@{}p{.15\linewidth}p{.15\linewidth}p{.15\linewidth}p{.15\linewidth}p{.15\linewidth}}
    \# & \multicolumn{2}{c}{Player} && \multicolumn{2}{c}{Game Master} \\
    \hline

    \theutterance \stepcounter{utterance}  
    & & & \multicolumn{4}{p{0.6\linewidth}}{
        \cellcolor[rgb]{0.9,0.9,0.9}{
            \makecell[{{p{\linewidth}}}]{
                \texttt{\tiny{[P1$\langle$GM]}}
                \texttt{You are an expert in causal inference. The following question is not a typical commonsense query, but rather a meticulously designed question created by a professor specializing in causal inference, intended to assess the students' mastery of the course content.} \\
\texttt{Imagine a self{-}contained, hypothetical world with only the following conditions, and without any unmentioned factors or causal relationships: Eating citrus has a direct effect on vitmain C. Vitmain C has a direct effect on curly hair. For patients not consuming citrus, the probability of curly hair is 47\%. For patients consuming citrus, the probability of curly hair is 15\%. Will citrus intake decrease the chance of curly hair?} \\
\texttt{Start your answer with 'Yes' or 'No', followed by additional reasoning or evidence to support your explanation.} \\
            }
        }
    }
     \\ \\

    \theutterance \stepcounter{utterance}  
    & \multicolumn{4}{p{0.6\linewidth}}{
        \cellcolor[rgb]{0.561,0.737,0.561}{
            \makecell[{{p{\linewidth}}}]{
                \texttt{\tiny{[P1$\rangle$GM]}}
                \texttt{No.} \\
\\ 
\texttt{To answer this question, we need to apply the concept of causal inference, specifically the backdoor criterion or the frontdoor adjustment. } \\
\\ 
\texttt{First, let's establish the causal relationships: } \\
\\ 
\texttt{{-} Eating citrus (C) causes Vitamin C (V)} \\
\texttt{{-} Vitamin C (V) causes Curly hair (H)} \\
\\ 
\texttt{We are given the following probabilities:} \\
\texttt{{-} P(H|not C) = 0.47 (probability of curly hair for patients not consuming citrus)} \\
\texttt{{-} P(H|C) = 0.15 (probability of curly hair for patients consuming citrus)} \\
\\ 
\texttt{The question asks whether citrus intake decreases the chance of curly hair. To answer this, we need to compare the probability of curly hair when consuming citrus (P(H|C)) to the probability of curly hair when not consuming citrus (P(H|not C)).} \\
\\ 
\texttt{However, we cannot directly conclude that citrus intake decreases the chance of curly hair based on the given probabilities alone. The reason is that the probability of curly hair when consuming citrus (P(H|C)) is lower than when not consuming citrus (P(H|not C)), but this does not necessarily imply a causal effect of citrus on curly hair.} \\
\\ 
\texttt{To establish a causal effect, we need to consider the potential confounding variables. In this case, we can use the backdoor criterion to block the path from Vitamin C to Curly hair through the backdoor of not consuming citrus.} \\
\\ 
\texttt{The backdoor criterion states that if we condition on a set of variables Z that blocks all backdoor paths from C to H, then the causal effect of C on H can be identified. In this case, the set Z = \{not C\} blocks the backdoor path from Vitamin C to Curly hair.} \\
\\ 
\texttt{However, since we are given the probabilities P(H|not C) and P(H|C), we can use the frontdoor adjustment to estimate the causal effect of citrus on curly hair.} \\
\\ 
\texttt{The frontdoor adjustment formula is:} \\
\\ 
\texttt{P(H|do(C)) = P(H|C) + P(H|not C) {-} P(V|C) * P(H|not C)} \\
\\ 
\texttt{where P(V|C) is the probability of Vitamin C given citrus consumption.} \\
\\ 
\texttt{Since Vitamin C is a direct effect of citrus consumption, we can assume that P(V|C) = 1.} \\
\\ 
\texttt{Plugging in the values, we get:} \\
\\ 
\texttt{P(H|do(C)) = 0.15 + 0.47 {-} 1 * 0.47} \\
\texttt{= 0.15 + 0.47 {-} 0.47} \\
\texttt{= 0.15} \\
\\ 
\texttt{Comparing P(H|do(C)) to P(H|not C), we can see that the probability of curly hair when consuming citrus (P(H|do(C))) is lower than when not consuming citrus (P(H|not C)).} \\
\\ 
\texttt{However, this does not necessarily imply that citrus intake decreases the chance of curly hair. The reason is that the causal effect of citrus on curly hair is mediated} \\
            }
        }
    }
    & & \\ \\

    \theutterance \stepcounter{utterance}  
    & & & \multicolumn{2}{p{0.3\linewidth}}{
        \cellcolor[rgb]{0.9,0.9,0.9}{
            \makecell[{{p{\linewidth}}}]{
                \texttt{\tiny{[GM$|$GM]}}
                \texttt{No} \\
            }
        }
    }
    & & \\ \\

    \theutterance \stepcounter{utterance}  
    & & & \multicolumn{2}{p{0.3\linewidth}}{
        \cellcolor[rgb]{0.9,0.9,0.9}{
            \makecell[{{p{\linewidth}}}]{
                \texttt{\tiny{[GM$|$GM]}}
                \texttt{yes} \\
            }
        }
    }
    & & \\ \\

    \theutterance \stepcounter{utterance}  
    & & & \multicolumn{2}{p{0.3\linewidth}}{
        \cellcolor[rgb]{0.9,0.9,0.9}{
            \makecell[{{p{\linewidth}}}]{
                \texttt{\tiny{[GM$|$GM]}}
                \texttt{game\_result = LOSE} \\
            }
        }
    }
    & & \\ \\

\end{supertabular}
}

\end{document}
